\subsection{An intrinsic finite quotient for the remaining arithmetic value}

\begin{statement}{B.39}{A class-two quotient of order \(3^{36}\) detects the remaining integral phase}{Proved}
Retain the actual local group $N$, the lattice
$\mathcal J=[\mathbb Z_3u+\mathbb Z_3v,\Lambda^2V](-1)$,
and the coefficient row $\lambda_{81}$ of B.38.
Then
\[
 \lambda_{81}(\operatorname{ad}_wV)=0,
 \qquad
 \lambda_{81}(\mathcal J)=3\mathbb Z/81.
                                                               \tag{B.39.1}
\]
Consequently division by three defines a primitive
local coefficient quotient
\[
 \lambda_{\mathcal J}:\mathcal J\twoheadrightarrow\mathbb Z/27.
                                                               \tag{B.39.2}
\]
The remaining cyclic difference in B.38 has value
$18\pmod{27}$ in its recorded orientation, and hence
is detected by this quotient.

There is an actual Frobenius-stable quotient $G$ of
the intrinsic geometric pro-three group $N$ with
\[
 |G|=3^{36},\qquad \exp G=27,\qquad
 G^{\rm ab}\simeq(\mathbb Z/27)^{10},\qquad
 G'\simeq(\mathbb Z/27)^2.
                                                               \tag{B.39.3}
\]
Its nilpotency class is two. The scalar
$\lambda_{\mathcal J}(\eta(F_2))$ is an explicit
primitive linear functional of the quadratic term
of the actual Frobenius automorphism of $G$.
This functional is unchanged by quadratic coordinate
changes and by changing the arithmetic Frobenius lift
by a geometric inner automorphism.

Thus the actual outer Frobenius action on the finite
group $G$ determines the remaining local phase.
The group, its stable kernel, and the required statistic
are specified below. B.40 evaluates that statistic as
$19\pmod{27}$ by the actual geometric extension.
The function-field equations of its cover and the full
quadratic Frobenius matrix are not required for that
evaluation and have not been computed here.
\end{statement}
\begin{proof}
\emph{A pair of alternating coefficient forms.}
Put $R=\mathbb Z/27$ and use the Hall coordinates of
B.35--B.38. Multiplication of the saved row
$\lambda_{81}$ by the ten primitive inner columns
gives zero modulo $81$. Multiplication by the ninety
primitive columns of $\mathcal J$ gives multiples of
three, including a multiple with valuation exactly
one. These exact calculations prove (B.39.1).
The division in (B.39.2) is the canonical isomorphism
$3\mathbb Z/81\simeq\mathbb Z/27$.
It therefore preserves local equivariance.

Define
\[
 \phi:\Lambda^2(V/27V)\longrightarrow R^2,
 \qquad
 \phi(b)=
 \begin{pmatrix}
 \lambda_{81}([u,b])/3\\
 \lambda_{81}([v,b])/3
 \end{pmatrix}.
                                                               \tag{B.39.4}
\]
The underlying row on the untwisted cubic lattice
has Frobenius multiplier two. Let $F$ be the actual
degree-one Frobenius matrix and $M_2$ its action on
the ordered pair $u,v$; thus
$F(u,v)=(u,v)M_2$ and $\det M_2=2$.
Equivariance gives
\[
 \phi\circ\Lambda^2F=C\phi,
 \qquad C=2M_2^{-t}.
                                                               \tag{B.39.5}
\]
In the fixed compatible quaternion basis,
\[
 M_2=\begin{pmatrix}1&23\\21&26\end{pmatrix},
 \qquad C=\begin{pmatrix}26&6\\4&1\end{pmatrix}
                  \pmod{27}.
                                                               \tag{B.39.6}
\]
The columns of $\phi$ at
$e_2\wedge e_0,e_1\wedge e_0$ form the minor
\[
                     \begin{pmatrix}25&13\\13&6\end{pmatrix},
 \qquad \det=-19\in R^\times.
                                                               \tag{B.39.7}
\]
Thus $\phi$ is onto $R^2$ and its kernel is a direct
summand of rank $43$.
All entries of (B.39.5), including the glued compact
and boundary blocks, are checked directly by
\path{c_residual_class_two_detector.py}.

\emph{The group quotient and its arithmetic stability.}
Class-two BCH is integral over $\mathbb Z_3$ because
$1/2\in\mathbb Z_3$. Hall collection gives
\[
 N/\gamma_3N\simeq V\oplus\Lambda^2V,
 \quad
 (x,z)(x',z')=(x+x',z+z'+\tfrac12x\wedge x').
                                                               \tag{B.39.8}
\]
A choice of free generators supplies the displayed
splitting; the central commutator subgroup itself is
canonical. The subgroup generated by twenty-seventh
powers is $27V\oplus27\Lambda^2V$, since
$(x,z)^{27}=(27x,27z)$.
After imposing these powers, quotient the central
commutator subgroup by $\ker\phi$.
This defines an actual quotient of $N$, with law
\[
 \begin{split}
 G&=(V/27V)\oplus R^2,\\
 (x,z)(x',z')&=
 (x+x',z+z'+14\phi(x\wedge x')).
 \end{split}                                                  \tag{B.39.9}
\]
Surjectivity of $N\to G$ follows from surjectivity on
the abelianization and from (B.39.7) on commutators.
Its kernel is the preimage of $\ker\phi$ together
with $\gamma_3N$ and the twenty-seventh powers.
The lower-central and power subgroups are
characteristic, and (B.39.5) makes the remaining
central kernel Frobenius-stable.
Good reduction with the fixed unit tangent makes
local inertia trivial. This proves stability under
the actual local arithmetic action, without knowing
its quadratic part.

Bilinearity proves associativity of (B.39.9), and
its commutator is $(0,\phi(x\wedge x'))$.
These identities and surjectivity of $\phi$ give
all assertions in (B.39.3).
Only the class-two BCH law has been used; no
class-three integral Lazard correspondence is assumed.

\emph{The precise Frobenius statistic.}
Every automorphism of (B.39.9) with the actual
graded actions has the form
\[
            \mathcal F(x,z)=(Fx,Cz+Bx),
       \qquad B\in\operatorname{Mat}_{2\times10}(R).
                                                               \tag{B.39.10}
\]
The extra term must be additive by (B.39.5), and
hence is $R$-linear. Conversely any such $B$ defines
an automorphism with these graded actions.
For the actual automorphism, $B$ is the projection
under $\phi$ of its integral quadratic action term.

Write
\[
 M_2=\begin{pmatrix}a&b\\c&d\end{pmatrix},
 \qquad B(u)=\binom{p_u}{q_u},
 \qquad B(v)=\binom{p_v}{q_v}.
\]
Use the integral pure-bracket lift of $w$ proved
in B.38. Its cubic variation depends only on the
quadratic action term $B_{F_2}:V\to\Lambda^2V$:
\[
 \eta(F_2)=\tfrac12
     \bigl([B_{F_2}u,Fv]+[Fu,B_{F_2}v]\bigr).
\]
Applying (B.39.2) gives the required statistic
\[
 \boxed{
 \lambda_{\mathcal J}(\eta(F_2))
     =\tfrac12(-bp_u-dq_u+ap_v+cq_v)\pmod{27}.}
                                                               \tag{B.39.11}
\]
The exact twenty-entry row on $B$ is recorded in
the script's JSON and NPZ files, and is primitive.
Although $\phi(w)=0$, formula (B.39.11) need not
vanish: it measures the arithmetic central correction
on lifts of the pair, rather than their commutator
inside the class-two quotient.

\emph{Both possible coordinate ambiguities disappear.}
Changing the splitting to $z'=z+Ex$ replaces $B$ by
$B+EF-CE$.
Put $E_H=E(u,v)$ and
$J_0=\left(\begin{smallmatrix}0&1\\-1&0\end{smallmatrix}\right)$.
For $B_H=B(u,v)$ the right side of (B.39.11) is
\[
               P(B_H)=2^{-1}
                   \operatorname{tr}((M_2J_0)^tB_H).
\]
The identities $M_2J_0M_2^t=2J_0$ and
$C=2M_2^{-t}$ give
\[
 P(E_HM_2)=\operatorname{tr}(J_0^tE_H)=P(CE_H).
\]
Thus every quadratic splitting change is killed.
Changing the arithmetic lift by an inner automorphism
of abelian image $t$ instead adds
$x\mapsto\phi(t\wedge Fx)$ to $B$.
The Jacobi identity identifies its contribution
with $\lambda_{\mathcal J}([t,w])$, which is zero
by (B.39.1). The script independently checks these
identities on all twenty coordinate changes and
all ten inner changes.

Finally B.38's oriented residual value $54$ becomes
$18$ after division by three. B.36 gives that the
actual value of (B.39.11) is one modulo three.
The graded matrices alone cannot choose it either:
(B.39.10) permits every linear $B$, and the statistic
is primitive. The theorem specifies the actual finite
group and the exact outer-arithmetic datum for this
test. The independent geometric-module calculation in
B.40 supplies its value $19$, without calculating the
full matrix $B$ or equations for the finite cover.
\end{proof}
